% ===================== ANEXOS =====================
\appendix
\section{Código de Simulación Hidráulica en Python}
\label{appendix:code}

Este anexo contiene el código fuente desarrollado en Python 3.10 para realizar la simulación hidráulica dinámica comparativa y el cálculo hidrostático estático en el fondo de la perforación dirigida. El código implementa las ecuaciones de Darcy-Weisbach y la aproximación de Haaland para resolver de forma explícita las pérdidas dinámicas.

\begin{verbatim}
import math

# Constantes físicas
g = 9.80665  # m/s^2

# Propiedades de los fluidos
fluidos = {
    "Salmuera Saturada (26% NaCl)": {
        "densidad": 1200.0,    # kg/m^3
        "viscosidad": 0.0016,  # Pa*s
        "flujo_masico": 300000.0  # kg/h
    },
    "Condensado (Agua Industrial)": {
        "densidad": 1000.0,    # kg/m^3
        "viscosidad": 0.0010,  # Pa*s
        "flujo_masico": 300000.0  # kg/h
    }
}

# Parámetros geométricos (HDPE SDR 17 - 12" NPS)
diametro_hdpe_12in = {
    "OD": 0.315,  # m
    "ID": 0.278   # m
}
rugosidad_hdpe = 1.5e-6  # m

# Longitudes del tramo
L_superficial = 450.0
L1 = math.sqrt(70.0**2 + (2558.35 - 2542.75)**2)
L2 = 180.0
L3 = math.sqrt(200.0**2 + (2557.56 - 2542.75)**2)
L_phd = L1 + L2 + L3

K_superficial = 0.2
K_phd = 0.5

def calcular_factor_friccion_darcy(Re, roughness_ratio):
    if Re < 2300:
        return 64.0 / Re
    else:
        termino = (roughness_ratio / 3.7) ** 1.11 + 6.9 / Re
        inv_sqrt_f = -1.8 * math.log10(termino)
        return 1.0 / (inv_sqrt_f ** 2)

def calcular_perdidas_tramo(Q_m3h, ID, L, K, rho, mu):
    Q = Q_m3h / 3600.0  # m^3/s
    A = (math.pi / 4.0) * (ID ** 2)
    v = Q / A
    Re = (rho * v * ID) / mu
    roughness_ratio = rugosidad_hdpe / ID
    f = calcular_factor_friccion_darcy(Re, roughness_ratio)
    h_f = f * (L / ID) * (v ** 2) / (2.0 * g)
    h_m = K * (v ** 2) / (2.0 * g)
    h_total = h_f + h_m
    dP_total = rho * g * h_total
    return v, Re, h_total, dP_total / 1000.0
\end{verbatim}

\section{Planos y Documentación Técnica de Referencia}
\label{appendix:drawings}

La Tabla~\ref{tab:documentos_referencia} relaciona los planos de diseño civil e hidráulico y la correspondencia técnica oficial consultada como base para el presente estudio técnico.

\begin{table}[htbp]
    \centering
    \caption{Documentación de referencia empleada en el estudio.}
    \label{tab:documentos_referencia}
    \setcounter{itemcount}{0}
    \begin{tabularx}{\textwidth}{@{}NllX@{}}
        \toprule
        \multicolumn{1}{c}{\textbf{Item}} & \textbf{Documento} & \textbf{Código / Código Rev.} & \textbf{Descripción} \\
        \midrule
        & Plano de Perfil PHD & `202609-BRI-SE-CIV-PL-001-RB` & Plano civil e hidráulico de planta y perfil longitudinal de la Perforación Horizontal Dirigida. \\
        & Correo de Solicitud BRINSA S.A. & Fwd: Validación DML (23/06/2026) & Correo electrónico oficial que especifica los escenarios operativos, caudales e interrogantes técnicos. \\
        \bottomrule
    \end{tabularx}
\end{table}
